
\section{Convergence and numerical controls}\label{rep:v-convergence}

This section collects the evidence that the numbers of Part~I are converged and that their numerical errors are small compared with the model items.
We find that the critical top mass has converged in the numerical implementation of the reduced model to $0.001\GeV$, that every numerical control of the Higgs chain together moves it by at most $0.002\GeV$, the ``numerics'' item of the error budget, and that the statistical and discretization errors of the shape and microwave-background calculations are below the precision at which their results are quoted. %
The section has three parts: the convergence of the critical mass through the physical refinements of the model (Sec.~\ref{sec:v20-refinements}), the numerical controls of each chain (Secs.~\ref{sec:v20-higgs}--\ref{sec:v20-cmb}), and the verification of the code package against the archived original runs (Sec.~\ref{sec:v20-package}).

\subsection{The critical mass through the refinements of the model}\label{sec:v20-refinements}

Table~\ref{tab:rb-convergence} of Sec.~\ref{rep:h-budget} follows the critical top mass at $\Pone=0.9$ and $1/2$ through the refinements of the Higgs zero-mode model, all with the same hazard procedure at $\alpha_s(M_Z)=0.1180$. %
The table is a record of how the model was built, not a sequence of successive orders of one expansion, and it should be read with that in mind.


The first five steps add distinct physical effects: the $\MSbar$ continuum term of the finite-volume matching and the Coulomb-gauge counting of the massive vectors (Sec.~\ref{rep:h-matching}), the exact gauge shape from DMC and the joint dynamics of the hypercharge mode (Sec.~\ref{rep:h-reduced}).
Each moves $m_c$ by $0.04$ to $0.11\GeV$ toward larger values. %
The next three rows test the numerical implementation of one and the same model, the two-dimensional model with the DMC shape.
It was first evaluated by converting widths with ratios computed at fixed couplings; a calibrated one-dimensional hazard and the direct two-dimensional roots at every node change $m_c$ by only $0.011$ and $0.001\GeV$. %
This is the sense in which the numerical implementation of the reduced model has converged: the remaining rows change the physics, not the numerics.
The NNLO mapping of Sec.~\ref{rep:h-smmap} shifts all levels by $+0.063\GeV$ through $\lambda(\mub;m_t)$ alone. %
The cross terms of Sec.~\ref{rep:h-cross} are the only refinement that lowers $m_c$; their final one-loop terms (background-covariant vev mixing, the Coulomb mismatch and the top chiral loop) move it back by $+0.005$ to $+0.007\GeV$. %
The vertex row uses the anticommutator vertex that the final calculation replaced; it is kept to show the size of the step, which is $-0.003\GeV$. %
The final critical masses are $172.502$, $172.960$ and $173.418\GeV$ at $\Pone=0.9$, $0.5$ and $0.1$; the package reproduces the full table from the archived summaries (\texttt{higgs/06\_hazard\_criterion}, key \texttt{convergence\_table\_0p1180}). %

\subsection{Numerical controls of the Higgs chain}\label{sec:v20-higgs}

\paragraph{Resonance roots.}
The two-dimensional Siegert roots are computed with a sine discrete-variable representation of $N$ points on $(0,S_{\max})$ in $s$, a renormalized matrix Numerov integration with step $h$ in $r$, and a second-order WKB outgoing condition at $R_{\rm out}$ (Sec.~\ref{rep:h-reduced}).
At $\lambda=-0.010$ and $-0.0143$ with the one-loop couplings, changing $(N,S_{\max})$ from $(32,14)$ to $(48,16)$ or $(24,12)$ changes $\Gamma$ by $2\times10^{-10}$; halving the step to $h=0.001$ changes it by at most $8\times10^{-8}$; moving $R_{\rm out}$ to $14$ by $7\times10^{-7}$; and dropping the second-order terms of the outgoing condition by $3.4\times10^{-3}$ (relative), which measures the order of the boundary condition rather than an error of the standard setting. %
The standard configuration of the hazard, $N=24$, $S_{\max}=12$, $h=0.004$, $R_{\rm out}=16$, differs from the reference configuration by at most $1.2\times10^{-6}$ in $\Gamma$; $R_{\rm out}=14$ changes $\Gamma$ by $2.7\times10^{-6}$ and a zeroth-order outgoing condition by $3.0\times10^{-3}$. %
On the critical mass these are below $4\times10^{-7}$, $9\times10^{-7}$ and $1.0\times10^{-3}\GeV$. %
Every root is recomputed with a larger $s$ box, and a change $|\Delta E|>\max(2\times10^{-7},\,2\times10^{-3}\Gamma)$ flags a box state; in the runs with the earlier mapping all 377 two-dimensional columns passed at the first attempt, with $|\Delta E|\le1.4\times10^{-11}$ and relative changes of $\Gamma$ of at most $7.8\times10^{-5}$, and the 220 columns with the NNLO mapping passed in the same way. %
Eleven nodes instead of 29 change $m_c$ by $7\times10^{-5}\GeV$. %
With the one-loop couplings the node code reproduces the static widths of the reduced-model calculation to $4.4\times10^{-6}$. %

\paragraph{Independent width methods.}
Three methods with different boundary treatments agree with the Siegert roots.
Finite differences with the complex absorbing potential $-i\eta(r-7)^2\theta(r-7)$ give $\Gamma_{\rm 2D}/\Gamma_{\rm 1D}=0.5313$ and $0.5317$ at $h=0.04$ and $0.03$, against $0.5303$ from the Siegert roots; the absolute widths of both models are low by $2.6\%$, which cancels in the ratio. %
A real wall moved from $R_w=8$ to $11$ at $\lambda=-0.0143$ shows one level that stays at $\varepsilon=4.2737\pm0.002$ while the others cross it, matching the Siegert value $4.27375$. %
For the one-dimensional model exterior complex scaling agrees with the Siegert root for the same inner boundary to $3\times10^{-9}$ in $E$ and $3\times10^{-7}$ relative in $\Gamma$. %
The absorbing-potential eigenvector has $P(r<5)=0.9988$ at $\lambda=-0.010$ and probability $10^{-10}$ beyond $0.75\,S_{\max}$ in $s$, so the resonance is localized in the well and the $s$ box plays no role. %

\paragraph{Analytic limits.}
A three-dimensional isotropic oscillator in the s wave gives $0.4500000000041$ against the exact $0.45$, and the uncoupled limit $\Omega\equiv0.3$ separates to $8\times10^{-11}$. %
Over $\lambda=-0.004$ to $-0.020$ the Kemble WKB width agrees with the exact width to $3\%$ for tunnelling exponents $S\approx2$--$25$. %
Changing the potential beyond $1.5$ times the outer turning point (linear tail, doubled quartic, frozen matching) changes $\Gamma$ by at most $0.75\%$, and beyond twice that radius by $10^{-4}$. %

\paragraph{Hazard integral and roots in $m_t$.}
The hazard $I=\int_0^5\Gamma\,d\tau/(hL_b\sqrt{\cosh2\tau})$ is evaluated from a cubic spline of $\ln\Gamma$ in $\tau$ on the eleven nodes and Simpson's rule on $20\,001$ points; the critical masses come from a cubic spline of $\ln I$ in $m_t$ on a $0.2\GeV$ grid and a Brent root.
An independent adaptive quadrature agrees with the standard integration to better than $10^{-12}\GeV$ in $m_c$.
Monotone PCHIP interpolants, which are less accurate than the cubic splines of the main calculation, move $m_c$ by at most $0.0002\GeV$ (interpolation in $\tau$) and $0.0025\GeV$ (in $m_t$); this is an upper bound on the interpolation sensitivity, compatible with the numerical item of the budget. %
For the background inputs the widths are computed at the eleven nodes and at three further scales and interpolated in $\ell=\ln(\mu/\mub)$; four direct reruns with all nodes shifted agree with the interpolation to $1.1\times10^{-7}\GeV$. %

\paragraph{Diffusion Monte Carlo.}
The zero-mode DMC code was validated on problems with known spectra: the quartic oscillator gives $2.1402(5)$ against the certified $2.1407$, and two further benchmark levels give $4.1171(14)$ and $12.5869(22)$ against $4.116720$ and $12.5868$. %
The walled DMC of Sec.~\ref{rep:v-fulldim} reproduces the unwalled value at $\lambda=0$ to $0.75\sigma$; its time-step extrapolation and statistics are described there.

\paragraph{Cross terms.}
The logarithms of the one-loop cross terms close the renormalization group of the Higgs field: the closure ratio is $0.994$, $0.973$, $0.993$ and $0.990$ at $r=2$, $3$, $4$ and $5$, so $\mu=1/L$ is a consistent matching scale (Sec.~\ref{rep:h-cross}). %
The top chiral logarithm satisfies $B/(q_\Phi^2\tilde m_t^2)=1.000\pm0.001$, and the groups of cross terms add to the total to $2\times10^{-4}\GeV$. %
A recomputation of all 22 node jobs of the final run from scratch in the package (701~s on eleven workers, gauge profiles recomputed) is bit-identical to the original shards, and so are the critical masses.

\paragraph{Standard Model mapping and matching-scale scan.}
The Python port of the SMDR matching and running was cross-checked against the compiled SMDR~1.3 at all matching scales, and the driver used for the compiled refits reproduces the archived fits bit for bit (Sec.~\ref{rep:h-smmap}).
The matching-scale scan of the final model reproduces the final roots and critical masses at $c=1$ to $2\times10^{-12}$. %

\paragraph{Marginalization.}
The quadrature of the marginalized probability and an independently written, seeded Monte Carlo integration with a different interpolating surface agree to $5\times10^{-4}$ for all eleven mass definitions; the curves at $\alpha_s\neq0.1180$, obtained from a quadratic map, agree with direct hazards at $\alpha_s=0.1171$ and $0.1189$ to $|\Delta\Pone|\le4.4\times10^{-4}$ (Sec.~\ref{rep:h-budget}). %

\subsection{Numerical controls of the shape chain}\label{sec:v20-shape}

\paragraph{Casimir sums and identities.}
The Casimir energies are exact image sums evaluated with an Ewald split at $s=2$ and checked against direct sums; the shape Hessians use five-point finite differences along geodesics $G=\exp(2\sum_A\beta_AX_A)$.
The trace identity ${\rm Tr}\,{\rm Hess}\,Z_{\rm eff}=\frac83Z_{\rm eff}$ holds to $10^{-8}$ for all eight spin structures (Sec.~\ref{rep:casimir}). %
The no-go identities of Sec.~\ref{rep:shape} were checked symbolically and with $40$--$60$-digit numerics (Chowla--Selberg and Ewald forms in mpmath), and the closed-form gain $\mathcal G(\kappa)=\sqrt2\cosh(\pi\sqrt{4\kappa-1}/4)$ against direct integration of the mode equation. %

\paragraph{Linear shape layer.}
The closed-form floor, $\sigma_{\rm floor}^2=K(\kappa_E)/X$ with $K(19.32)=24.550$, the correlation lemma behind the real-moment theorem and the tables of $G(|t_x|)$ were checked against direct integration and in Wolfram Language (Sec.~\ref{rep:quantum-shape}); an independent code reproduced $G(|t_x|)$ (Sec.~\ref{rep:v-blind}). %

\paragraph{Nonlinear trajectories.}
The nonlinear model evolves the volume and the five shape moduli with the twisted Ewald force of the Standard Model and the zero-mode energy, with the Dormand--Prince integrator DOP853 at relative tolerance $10^{-11}$ and absolute tolerance $10^{-13}$; a trajectory is declared pinched when the smallest eigenvalue of the shape metric falls below $10^{-3}$. %
(A header comment of the original script stated a pinch threshold of $10^{-6}$; the original runs used $10^{-3}$, and an erratum is archived with the code.)
Initial conditions are sampled from the truncated Wigner function with recorded seeds and mirror pairing, which makes the two-branch survival exactly symmetric under $t_x\to-t_x$; re-integrating archived trajectories reproduces them bit for bit.
Survival fractions are binomial and are quoted with Wilson $95\%$ intervals.
The levels of the conditional theorem were recomputed with $N=2048$ paired draws per node: $t_*(99\%)=0.237/h$ with the $95\%$ range $0.225$--$0.252/h$, and $t_*(90\%)\approx0.35/h$ (logistic fit $0.352$; $0.355$ from interpolation between nodes), so the quoted ``$0.35$ and $0.24$ Hubble times'' carry their statistical error. %
For the amplitude threshold the log-linear and logistic fits give $s_{99}=2.76$ and $2.86$ at $N=2048$ and an independent code gives $2.95$, which is why the papers quote $s_{99}\approx2.8$--$2.95$. %
On the $A_s$-calibrated plateau history instead of the leading $\cosh$ background none of $2560$ paired trajectories changes its fate. %

\paragraph{Quantum corrections to the sampling.}
The leading Moyal correction to the fossil variance is $-2.03\times10^{-4}$ (aspect) and $-1.08\times10^{-3}$ (angle) relative; in a one-mode model the exact quantum evolution and the truncated Wigner evolution differ by $-0.011\%$ (Sec.~\ref{rep:quantum-shape}). %

\paragraph{Emptiness gate and primordial spectra.}
The mode integrations of the emptiness gate conserve the Wronskian to $4\times10^{-10}$ (standard with 97 and refined with 401 wave numbers). %
For the primordial spectra, changing the initial state (adiabatic order $0$ or $2$, $6$--$10$ e-folds of preparation) changes $n_s$ by less than $10^{-7}$ (Sec.~\ref{rep:spectra}). %

\subsection{Numerical controls of the microwave-background chain}\label{sec:v20-cmb}

The covariance of the temperature and $E$-mode fields on the torus is a sum over the discrete wave vectors with CAMB~2.0.4 transfer functions (with a fix of the low-$k$ continuation and a scan of the transfer cutoff), completed by the isotropic continuum above the cutoff (Sec.~\ref{rep:sky}).
The ensemble uses $x_{\max}=k_{\max}\chi_*=240$ with a smooth taper of width $0.25$.
At the representative edge, raising $x_{\max}$ to $320$ changes the divergence $D$ for $\ell\le80$ by $-0.7\%$, and halving the taper changes it by $+0.17\%$, which moves the read edge by about $10^{-4}\,\dLSS$. %
The Kullback--Leibler divergences use traces and log-determinants of packed inverses, with a relative error of at most $1.7\times10^{-12}$ in the eigenvalue form. %
The profile over the volume of the reference cube is a minimum over a pool of 19 computed cubes; it replaced an earlier parabolic profile that underestimated 82 pairs by more than $1\%$ and violated the data-processing inequality once, and the corrected table reproduces the peak readability $0.34398$ at $L_0=1.075\,\dLSS$ of the ensemble analysis bit for bit. %
Convergence was tested directly only at the representative edge; all readability statements are for the ideal full sky, $\ell\le80$, and the read criterion $D\ge4.5$. %

\subsection{Verification of the code package}\label{sec:v20-package}

Every result directory of the code package (Sec.~\ref{rep:v-repro}) has a list of checks with tolerances fixed in advance: critical top masses to $0.001\GeV$, trajectory records to $10^{-9}$ relative, and most analysis layers bit for bit.
A quick mode recomputes a non-trivial part of each calculation (a two-dimensional node job, a DMC shard, the hazard integral, the claims layer, re-integrated nonlinear trajectories, covariance builds) and compares it with the archived outputs; a full mode repeats all original runs.
Table~\ref{tab:v22-verify} of Sec.~\ref{rep:v-repro} lists the state of all checks.
Recomputing every covariance matrix of the microwave-background ensemble for $\ell\le80$ from the shipped transfer table (about 70~min on 12 workers, with a 15.5~GB cache) reproduces the tabulated divergences to $2.3\times10^{-13}$, and the window edges, the peak of $P_{\rm joint}$ and the integrated measure to all quoted digits (Sec.~\ref{sec:v22-verify}). %
Two optional comparisons of the Fisher chain need inputs that are not shipped with the package, so its full mode covers only the C59 part.
